\documentclass{tietreport}

% Import images
\usepackage{graphicx}

% Bibliography configuration
\usepackage[
  date=year,
  style=alphabetic,
  backend=biber,
  sorting=nty,
  maxnames=5,
  minnames=3,
  maxalphanames=4,
  minalphanames=3,
  backref=true,
  doi=false,
  isbn=false,
  url=false,
  eprint=false
]{biblatex}
\addbibresource{references.bib}

%% ==================================================
%% CUSTOMIZE THESE FIELDS FOR YOUR PROJECT
%% ==================================================

% Course/Lab header
\TitlePageHeader{%
  UCS 503- Software Engineering%
}

% Main project title
\ProjectTitle{%
  CPRecall%
}

% Subtitle or project description
\ProjectSubTitle{%
  A Personalized LeetCode Learning, Mastery, and
  Spaced-Recall System -- Progress Report%
}

% Student information (up to 4 students)
\author{%
  \texttt{1024030665} Subham Kumar \\
  \texttt{1024030666} Tanush Grover\\
  \texttt{1024030669} Shubhi Singh%
}

% Group number, degree, year, program
\TitlePageSubText{%
  Group: Nexus 3 BE Third Year, CSE%
}

% Faculty advisor/guide name
\AdvisorName{%
  Raghav B. Venkataramaiyer%
}

% Evaluation period and department info
\TitlePageFooterText{{%
    \large Mid-Semester Evaluation \\[0.2cm]
    \bfseries Computer Science and Engineering
    Department%
  } \\[0.15cm]
  Thapar Institute of Engineering and Technology,
  Patiala%
}

%% ==================================================

\begin{document}

% Generate title page
\maketitle

% Table of contents (auto-generated from chapters)
\tableofcontents

% ===== CHAPTER 1: INTRODUCTION =====

\chapter{Introduction}

\section{Background and Context}

Competitive programmers and interview
aspirants typically solve a large number of
LeetCode problems, but plain problem-solving
volume does not by itself guarantee long-term
retention of the underlying Data Structures
and Algorithms (DSA) patterns. A user may
solve a problem once, yet forget the pattern
or the optimal approach weeks later because
there is no structured mechanism that
(a) evaluates how well the problem was
actually solved, and (b) schedules a future
review at the right time.

CPRecal (previously named CodeRecall) is
being built to close this gap. It is a
personalized LeetCode learning, mastery, and
recall system that observes a user's solving
behaviour on LeetCode, analyses the submitted
code for its actual and optimal complexity,
computes a quantitative mastery score for the
problem, aggregates that mastery at the
DSA-pattern level, and finally schedules
future revisions using the FSRS (Free Spaced
Repetition Scheduler) algorithm~\cite{fsrs}.

\subsection{Problem Statement}

Given that a user solves problems on LeetCode
without any external tracking, there is
currently no automated way to:
\begin{itemize}
  \item objectively score how efficiently a
    specific problem was solved (in terms of
    time complexity, number of attempts, time
    taken, and hints used),
  \item separate a user's mastery of an
    individual problem from their mastery of
    the broader DSA pattern that problem
    belongs to, and
  \item decide \emph{when} that user should be
    made to revisit the problem or pattern so
    that retention is maximised.
\end{itemize}
CPRecal addresses this by combining a browser
extension for observation, an LLM-based code
analysis step, a backend scoring and
scheduling engine, and a dashboard for
review.

\section{Project Objectives}

\begin{enumerate}
  \item \textbf{Objective 1:} Capture a
    user's LeetCode solving activity (problem
    metadata, session time, submission
    history, and the final accepted code)
    automatically through a Tampermonkey
    browser extension, with no manual data
    entry beyond an optional hints-used
    count.
  \item \textbf{Objective 2:} Analyse each
    accepted solution's actual and optimal
    time/space complexity using an LLM, and
    convert this, along with the other raw
    solving signals, into a weighted
    \emph{Problem Mastery} score.
  \item \textbf{Objective 3:} Aggregate
    problem-level mastery into a
    \emph{Pattern Mastery} score per DSA
    pattern, and drive a spaced-repetition
    review schedule for each user using the
    FSRS algorithm, surfaced through a React
    dashboard.
\end{enumerate}

These objectives were defined and finalised
during the design phase reported in this
document; implementation proceeds directly
against this specification.

% ===== CHAPTER 2: METHODOLOGY & DESIGN =====

\chapter{Methodology and Design}

\section{System Architecture}

The overall data flow that has been finalised
for CPRecal is:
\[
\text{LeetCode} \rightarrow
\text{Tampermonkey Extension} \rightarrow
\text{LLM Code Analysis} \rightarrow
\text{Extension UI} \rightarrow
\text{CPRecal Backend} \rightarrow
\text{MongoDB} \rightarrow
\text{Problem/Pattern Mastery + FSRS}
\rightarrow \text{React Frontend}
\]

\subsection{High-Level Design}

The end-to-end solving flow that has been
designed and agreed upon is as follows:
\begin{enumerate}
  \item The user opens a LeetCode problem;
    the extension identifies the problem ID,
    title, titleSlug, difficulty, and
    tags/patterns.
  \item The extension tracks the solving
    session and the total time taken, and
    monitors submission status/history as the
    user attempts the problem.
  \item On an Accepted submission, the
    extension captures the accepted code,
    language, runtime, memory, submission ID,
    and timestamp.
  \item The accepted code and question name
    are sent to an LLM, which returns the
    user's actual time/space complexity, the
    optimal time/space complexity, and an
    explanation.
  \item The extension computes a
    \emph{Solution Efficiency} score from the
    actual-versus-optimal complexity and
    displays it, along with the rest of the
    solving data, in a floating widget. The
    user may additionally record hints used.
  \item The extension packages the raw
    solving data and the LLM analysis and
    submits it to the authenticated CPRecal
    backend.
  \item The backend computes the final
    Problem Mastery score, persists it along
    with the supporting raw data, converts it
    into an FSRS rating, and lets FSRS
    determine the next review date, storing
    the resulting FSRS state.
  \item The React frontend fetches the
    personalized dashboard, problem library,
    pattern dashboard, and problem-detail
    data from the backend.
\end{enumerate}

A clear responsibility split between the
three layers has also been finalised:
the Tampermonkey extension is a
\emph{collector and display layer} only; the
LLM is responsible purely for analysing code
complexity and does not decide mastery; and
the backend remains the single source of
truth for authentication, persisted mastery,
pattern aggregation, and FSRS review state.

\subsection{Technical Stack}

\begin{itemize}
  \item \textbf{Framework/Language:} React
    for the frontend~\cite{reactjs}, Node.js
    with Express for the backend
    API~\cite{expressjs}, and a Tampermonkey
    user script for the browser extension
    layer~\cite{tampermonkey}.
  \item \textbf{Database:} MongoDB, accessed
    through Mongoose object
    models~\cite{mongodb, mongoose}.
  \item \textbf{Scheduling Algorithm:} FSRS
    (Free Spaced Repetition Scheduler) for
    computing review intervals from mastery
    ratings~\cite{fsrs}.
  \item \textbf{External Integration:} An LLM
    API for accepted-code complexity
    analysis, invoked either directly from the
    extension or, if security/reliability
    considerations require it, relayed through
    the backend.
\end{itemize}

\section{Design Completed So Far}

\subsection{Repository and Module Structure}

A monorepo layout has been finalised to keep
the React client, Express server, shared
contracts, and Tampermonkey extension
separated while remaining easy to scale:
\texttt{client/} (pages, components,
services, hooks), \texttt{server/} (models,
controllers, routes, middleware, and a
\texttt{services/} layer further split into
\texttt{mastery/}, \texttt{fsrs/},
\texttt{llm/}, and \texttt{ingestion/}
sub-modules), \texttt{extension/}
(collectors, GraphQL queries, analysis,
widget UI, and storage), \texttt{shared/}
(constants, schemas, types), plus top-level
\texttt{scripts/}, \texttt{tests/}, and
\texttt{docs/} directories. Folder-level
responsibilities -- for example, that
\texttt{server/models} holds only Mongoose
persistence definitions while
\texttt{server/controllers} handles only
HTTP-level request/response logic -- have
been written down so that implementation
work can proceed consistently across the
team.

\subsection{Data Model}

The database design distinguishes shared,
global data from user-specific data. A single
shared \texttt{Problem} collection stores
\texttt{problemId}, title, \texttt{titleSlug},
difficulty, and patterns/tags. All
user-specific activity is represented through
separate collections --
\texttt{SolvingSession}, \texttt{Submission},
\texttt{AcceptedSolution},
\texttt{ProblemMastery},
\texttt{PatternMastery}, \texttt{FSRSState},
and \texttt{ReviewHistory} -- each referencing
a \texttt{userId}. Raw observations (sessions,
submissions, accepted code) are retained
rather than discarded, so that mastery can be
recalculated later if the scoring formula
changes.

\subsection{Problem Mastery Formula}

Problem mastery answers ``how well did this
specific user perform on this specific
problem?'' and has been assigned the
following weights:
\begin{itemize}
  \item Solution Efficiency -- 25\%
  \item Time Complexity -- 25\%
  \item Number of Submissions -- 20\%
  \item Time Taken (benchmarked against
    problem difficulty) -- 15\%
  \item Hints Used -- 15\%
\end{itemize}
The exact normalisation inside each component
is left open for tuning during
implementation, but the relative weighting
has been fixed.

\subsection{Problem Mastery vs. Pattern Mastery}

A clear conceptual separation has been made
between problem mastery (one user's score for
one problem) and pattern mastery (that user's
cumulative score for a DSA pattern, computed
from their performance across all problems
tagged with that pattern). Because a problem
can carry multiple pattern tags, the Problem
Details view is specified to show, for each
pattern, the user's cumulative pattern
mastery rather than the problem's own score --
keeping the two metrics visually and
conceptually distinct.

\subsection{FSRS Integration}

A mapping from the 0--100 Problem Mastery
score to the four FSRS review ratings has
been finalised:
\begin{itemize}
  \item 0--39 mastery $\rightarrow$ rating 1
  \item 40--59 mastery $\rightarrow$ rating 2
  \item 60--79 mastery $\rightarrow$ rating 3
  \item 80--100 mastery $\rightarrow$ rating 4
\end{itemize}
This rating, combined with the user's
existing FSRS review state and history, is
what FSRS~\cite{fsrs} uses to compute the
next review date. Due items are designed to
surface on the Home Dashboard's revision
flow.

\subsection{Frontend Views Specified}

The required frontend surfaces have been
specified in detail:
\begin{itemize}
  \item \textbf{Home Dashboard:} total solved
    problems, problems due today, patterns
    attempted, and a ``Today's Revision''
    list showing problem name, tags/patterns,
    cumulative pattern mastery, and a
    done/not-done review status.
  \item \textbf{Problem Library:} the full
    problem set with search and filters (e.g.
    pattern, status), keeping global problem
    data separate from user-specific solved
    status/mastery.
  \item \textbf{Pattern Dashboard:} a user's
    cumulative mastery broken down by DSA
    pattern, to highlight strong and weak
    areas.
  \item \textbf{Problem Details:} problem
    metadata, the user's individual problem
    mastery, solving statistics (time taken,
    submission count, runtime, memory), and
    per-pattern cumulative mastery shown on
    hover.
\end{itemize}

\subsection{Extension Widget Specification}

The floating in-page widget shown on LeetCode
has been specified to display: Problem,
Difficulty, tags, Time Taken, Submissions
Until Success, Runtime, Memory, Efficiency
Score, Your Complexity, Optimal Complexity,
Space, Hints Used, a ``Submit to CPRecal''
action, and an ``Open CPRecal Dashboard''
link. A visual mock-up of this widget has
been reviewed by the team as a reference for
implementation.

\subsection{Authentication and Data Isolation}

CPRecal is designed as a standard
authenticated application in which every
user-specific record (\texttt{SolvingSession},
\texttt{Submission}, \texttt{AcceptedSolution},
\texttt{ProblemMastery},
\texttt{PatternMastery}, \texttt{FSRSState},
\texttt{ReviewHistory}) must belong to exactly
one CPRecal user, while shared problem
metadata remains global. It has been
explicitly decided that backend-side
authorization -- not the frontend or the
extension -- must determine the authenticated
user for every private read or write, so that
one user can never access another user's
solving data.

% ===== CHAPTER 3: RESULTS & EVALUATION =====

\chapter{Results and Evaluation}

\section{Current Status}

At this stage of the project, the work
completed is the requirements analysis and
system design described in Chapter~2: the
end-to-end data flow, the repository
structure, the database schema, the Problem
Mastery and Pattern Mastery formulas, the
FSRS rating mapping, and the specification of
every frontend view and the extension widget.
These have been documented as stable
references (\texttt{ARCHITECTURE.md} and
\texttt{PROJECT\_CONTEXT.md}) so that
implementation can proceed consistently
across the team and be handed off to, or
picked up by, any contributor without
ambiguity.

\section{Design Verification}

Since implementation and formal testing have
not yet begun, verification at this stage has
consisted of walking through the 17-step
solving flow (Section~2.1.1) end-to-end on
paper to confirm that every raw data point
needed by the Problem Mastery formula, the
Pattern Mastery aggregation, and the FSRS
rating conversion is in fact collected
somewhere upstream, and that the
responsibility split between the extension,
the LLM, and the backend leaves no step
without an owner. This walkthrough did not
surface any missing data dependency in the
current design.

\section{Pending Work}

Coding of the individual modules (extension
collectors, backend controllers/services, and
the React frontend), along with unit,
integration, and performance testing, has not
yet started and will follow directly from the
design finalised in Chapter~2.

% ===== CHAPTER 4: CONCLUSION =====

\chapter{Conclusion}

\section{Summary of Work}

At this stage, the CPRecal project has moved
from an initial idea to a fully specified
system design. The three objectives set out
in Section~1.2 -- automated data capture,
LLM-based complexity analysis feeding into a
weighted mastery score, and FSRS-driven
spaced revision -- have each been translated
into a concrete, documented design: a
monorepo structure, a MongoDB schema
separating shared and user-specific data, an
explicit Problem Mastery formula, a clear
Problem-Mastery-versus-Pattern-Mastery
distinction, an FSRS rating mapping, and a
full specification of every frontend view and
the extension widget.

\section{Key Findings}

\begin{itemize}
  \item \textbf{Finding 1:} Separating
    \emph{problem} mastery from \emph{pattern}
    mastery early in the design avoids
    conflating a single lucky/unlucky attempt
    with a user's broader command of a DSA
    pattern, and this distinction has been
    threaded consistently through the
    database schema, the mastery formulas, and
    every frontend view.
  \item \textbf{Finding 2:} Keeping the
    Tampermonkey extension a pure
    collector/display layer, with the backend
    as the sole source of truth for mastery
    and review state, keeps the architecture
    flexible -- for instance, the LLM call can
    be moved from the extension to the backend
    later (for API-key or reliability reasons)
    without changing how mastery or FSRS
    scheduling work.
  \item \textbf{Finding 3:} Retaining raw
    observations (sessions, submissions,
    accepted code) rather than only the
    derived mastery score means the scoring
    formula can be revised later and applied
    retroactively, which was an important
    design decision made before any
    implementation began.
\end{itemize}

\section{Future Work and Recommendations}

\begin{enumerate}
  \item Implement the Tampermonkey collectors
    and the extension widget per the
    specification in Section~2.2.6.
  \item Implement the backend models,
    controllers, and the
    mastery/FSRS/ingestion services per the
    module structure in Section~2.2.1.
  \item Implement the React dashboard,
    problem library, pattern dashboard, and
    problem-details pages per
    Section~2.2.5~\cite{reactjs}.
  \item Finalise the exact normalisation
    formulas within each Problem Mastery
    component (Section~2.2.3).
  \item Begin unit and integration testing
    once the above modules are implemented.
\end{enumerate}

\section{Final Remarks}

The design phase reported here was necessary
before writing implementation code, since it
fixes the data contracts between the
extension, the LLM, and the backend, and
avoids rework once coding begins. The next
report will cover the implementation
undertaken against this design.

% ===== BIBLIOGRAPHY =====

% Print bibliography with heading in TOC
\printbibliography[heading=bibintoc]

\end{document}
